
\subsection{Evaluation Method}

Ablation-style comparison?

Rough idea:
\begin{itemize}
    \item Baseline condition: An agent operating from a weakly structured textual state summary or conventional data API response, without explicit HTML hypermedia controls for the current valid actions.
    \item Hypermedia-constrained condition: An agent operating strictly within the affordances provided by server-returned HTML hypermedia representations, rendering or selecting UI actions mapped to the current application state and valid controls.
\end{itemize}

\todo{Draw a figure to show the testing part <= Before that make it sure that this works, need more thought}

\subsection{Metrics we Concern}

Measurement Variables:
\begin{enumerate}
    \item \textbf{Invalid Action Suggestion Count:} Number of actions suggested or rendered by the agent that are not valid in the current backend state.
    \item \textbf{Action-Affordance Alignment:} Proportion of generated UI actions that correspond to affordances exposed in the current hypermedia representation.
    \item \textbf{State Transition Correctness:} Whether executing a selected action leads to the expected next state and refreshed affordance set.
    \item \textbf{Interaction Step Count (maybe):} Total number of user-agent interactions required to reach task completion.
\end{enumerate}

\subsection{Experiment Results}

\todo{Core part of this chapter, may divide into more sections

- Screenshots or narative logs of the exp, give a direct feeling of how it works

- Tables and Figures, Visualization of metrics
}
